\documentclass[11pt,a4paper]{article}

\usepackage[margin=20mm]{geometry}
\usepackage[T1]{fontenc}
\usepackage{lmodern}
\usepackage{tabularx,booktabs,longtable,array}
\usepackage{hyperref}
\usepackage{xcolor}
\usepackage{enumitem}
\usepackage{graphicx}
\usepackage{float}

\hypersetup{
    colorlinks=true,
    linkcolor=blue,
    urlcolor=blue,
    citecolor=blue
}

\setlength{\parindent}{0pt}
\setlength{\parskip}{5pt}
\renewcommand{\arraystretch}{1.15}

\newcommand{\verify}[1]{\textcolor{red}{[VERIFY: #1]}}
\newcommand{\code}[1]{\texttt{#1}}

\title{\textbf{23CSE455 Design Patterns}\\
Case Study Report}
\date{28 October 2026}

\begin{document}
\maketitle

\begin{center}
\textbf{Case Study: Design Pattern Discovery and Evolution
in a Django Job Portal}
\end{center}

\textbf{Team:} B3

\textbf{Application:} Django Job Portal

\textbf{Submission repository:}
\url{https://github.com/aswathymohan-amrita/23CSE455_DP_CaseStudy}

\textbf{Students:}

\begin{tabularx}{\linewidth}{@{}lXl@{}}
\toprule
Roll Number & Student Name & Batch\\
\midrule
AM.SC.U4CSE23233 & M Poojith & B\\
AM.SC.U4CSE23234 & M Surya Prakash & B\\
AM.SC.U4CSE23241 & P Rushikesh & B\\
AM.SC.U4CSE23247 & S Krishna Vamsi & B\\
\bottomrule
\end{tabularx}

% ============================================================
\section{Application and Architecture}
% ============================================================

\subsection{Application Overview}

This case study examines a Django Job Portal application
that supports job seekers and employers through REST and
GraphQL interfaces. Its functionality includes job creation,
job applications, company management, applicant management,
permissions, and serialization.

The inspected implementation uses Django REST Framework
(DRF) for REST APIs and Graphene for GraphQL mutations.
Part 1 identifies and evaluates pattern collaborations in
the original application. Part 2 examines three changes
that introduce or extend design structures in the application.

\subsection{Source and Environment}

\begin{itemize}[leftmargin=*]
    \item Repository:
    \url{https://github.com/manjurulhoque/django-job-portal.git}

    \item Source revision: baseline commit
    \code{f6e5efb33e9aac7334a364ee46d6e2ad1321db30}.

    \item Access date: 28 October 2026. The repository's license
    and exact dependency versions should be cited from the checked-out
    repository and requirements files if required by the submission.

    \item Frameworks: Django, Django REST Framework, and Graphene;
    exact package versions are not reproduced here because the available
    evidence does not establish a complete version manifest.

    \item Database: use the database configured by the project settings;
    the baseline test log reports a Django in-memory SQLite test database.
\end{itemize}

\subsection{Scope and Method}

The baseline inventory documents five pattern instances
across the inspected GraphQL and REST API code. These
instances comprise two Strategy instances, two Template
Method instances, and one qualified Factory Method candidate.

The principal inspected files are:

\begin{itemize}[leftmargin=*]
    \item \code{jobsapp/graphql/permissions.py}
    \item \code{jobsapp/graphql/graphql\_mixins.py}
    \item \code{jobsapp/graphql/sub\_mutations.py}
    \item \code{jobsapp/api/permissions.py}
    \item \code{jobsapp/api/views/employer.py}
    \item \code{jobsapp/api/views/employee.py}
    \item \code{jobsapp/api/serializers.py}
\end{itemize}

Candidates were evaluated by examining their intent,
participant roles, inheritance, delegation, method calls,
and execution flow. Framework-provided mechanisms were
distinguished from application-defined behavior.

The findings describe the inspected scope and do not claim
that every file in the entire repository was exhaustively
audited.

\subsection{Architecture}


The application has two principal request paths.

\textbf{REST path.}
An HTTP request is dispatched to a DRF view. The configured
permission classes determine whether the request is allowed.
The view validates and processes data through serializers
and interacts with application models.
\begin{figure}[h]
    \centering
    \includegraphics[width=0.7\textwidth]{arch.png}
    \caption{System Architecture Diagram}
    \label{fig:architecture}
\end{figure}

\textbf{GraphQL path.}
A GraphQL mutation is dispatched through Graphene. Mutation
classes use permission mixins and operation-specific logic.
The common mutation workflow checks permissions before
delegating to the operation implementation.

Both paths interact with the application's models and
persistence layer.


\textbf{Architecture summary.} The REST path dispatches requests to DRF
views, checks configured permission classes, validates input with
serializers, and accesses application models. The GraphQL path dispatches
queries and mutations through Graphene; mutation mixins perform permission
checks and delegate operation-specific work. Both paths use the same
application persistence layer. A separate architecture image can be added
when the final image file is placed beside this source file.

\subsection{Run Instructions and Test Evidence}

Use the project's README and dependency files for setup,
database configuration, migrations, application startup,
and testing.

The recorded baseline command was \code{python manage.py test --keepdb}.
Django discovered 59 tests. Of these, 41 were executed: 31 passed,
5 failed, and 5 errored. A further 18 tests were not executed because
four test classes failed during \code{setUpClass}. This is a mixed
baseline result, not a fully passing test suite. A separate Part 1
pattern-test suite was reported as 17 passing tests; it is a distinct run
and is not added to the baseline total.

For reproducibility, run the project from its documented virtual
environment, install the dependencies in the repository's requirements
file, configure the database and environment variables, then run
the command above. Exact package versions and a complete setup transcript
are not established by the test logs available for this report.
% ============================================================
\section{Part 1: Pattern Discovery}
% ============================================================

\subsection{Discovery Method}

Each candidate was evaluated using the following criteria:

\begin{itemize}[leftmargin=*]
    \item The design problem addressed by the collaboration.
    \item The responsibilities of the participating classes
    and functions.
    \item Inheritance, composition, and delegation relationships.
    \item The execution flow between relevant methods.
    \item Whether the pattern belongs to application code
    or is supplied by a framework.
\end{itemize}

A class name or framework annotation alone was not
considered sufficient evidence of a design pattern.

\subsection{Baseline Totals}

\begin{center}
\begin{tabularx}{0.85\linewidth}{@{}Xr@{}}
\toprule
Measure & Count\\
\midrule
Documented pattern types & 3\\
Documented pattern instances & 5\\
Strategy instances & 2\\
Template Method instances & 2\\
Factory Method candidates & 1\\
\bottomrule
\end{tabularx}
\end{center}

The final totals for concrete participant classes, unique
participant symbols, and rejected and uncertain candidates
must be derived from the complete source and candidate
records. They are not inferred from the five inventory rows.

\subsection{Pattern Inventory}

\begin{longtable}{@{}p{0.20\linewidth}p{0.37\linewidth}p{0.35\linewidth}@{}}
\toprule
Pattern type & Instance & Main participants\\
\midrule
Strategy & GraphQL Permission Policy &
BasePermission, PermissionMixin, concrete policies\\

Strategy & DRF Role Permission Policy &
DRF BasePermission, IsEmployer, IsEmployee, IsJobCreator\\

Template Method & GraphQL Mutation Pipeline &
MutationMixin.mutate, CreateNewJobMixin, UpdateJobMixin\\

Template Method & DRF Create Hook &
CreateModelMixin.create, perform\_create overrides\\

Factory Method candidate & Company Serializer Selection &
get\_serializer\_class, CompanySerializer,
CompanyWriteSerializer\\
\bottomrule
\end{longtable}

\subsection{Instance 1: GraphQL Permission Policy}

\textbf{Pattern:} Strategy

\textbf{Category:} Behavioral

\textbf{Origin:} Application-defined implementation
\begin{figure}[h]
    \centering
    \includegraphics[width=0.7\textwidth]{uml1.png}
    \caption{Instance 1: GraphQL Permission Policy}
    \label{fig:architecture}
\end{figure}

\textbf{Participants:}

\begin{itemize}[leftmargin=*]
    \item Strategy interface: \code{BasePermission},
    in \code{jobsapp/graphql/permissions.py:9--16}.
    \item Concrete strategies: \code{AllowAny},
    \code{IsAuthenticated}, \code{IsEmployer}, and
    \code{IsEmployee}, in lines 23--46.
    \item Context: \code{PermissionMixin}, particularly
    \code{get\_permissions()} and \code{check\_permissions()},
    in \code{jobsapp/graphql/graphql\_mixins.py:17--49}.
    \item Clients: \code{CreateNewJob} and \code{UpdateJob},
    in \code{jobsapp/graphql/sub\_mutations.py:18--56}.
\end{itemize}

\textbf{Intent and evidence:}

GraphQL mutations require different access rules.
\code{BasePermission} defines common permission methods,
while concrete classes implement specific access decisions.

Mutation classes configure their \code{permission\_classes}.
The context instantiates the configured classes and calls
their permission methods. All configured policies must pass
for the permission check to succeed.

The collaboration is classified as Strategy because the
permission decision is delegated to configurable concrete
policies rather than being embedded in each mutation.


\subsection{Instance 2: GraphQL Mutation Pipeline}

\textbf{Pattern:} Template Method

\textbf{Category:} Behavioral

\textbf{Origin:} Application-defined workflow using mixins
\begin{figure}[h]
    \centering
    \includegraphics[width=0.7\textwidth]{uml2.png}
    \caption{Instance 2: GraphQL Mutation Pipeline}
    \label{fig:architecture}
\end{figure}

\textbf{Participants:}

\begin{itemize}[leftmargin=*]
    \item Template method: \code{MutationMixin.mutate()},
    in \code{jobsapp/graphql/graphql\_mixins.py:57--72}.
    \item Create hook: \code{CreateNewJobMixin.resolve\_mutation()},
    in lines 303--318.
    \item Update hook: \code{UpdateJobMixin.resolve\_mutation()},
    in lines 321--335.
    \item Concrete mutations: \code{CreateNewJob} and
    \code{UpdateJob}, in \code{sub\_mutations.py:18--56}.
\end{itemize}

\textbf{Intent and evidence:}

The common mutation workflow checks permissions before
delegating to operation-specific behavior. The
\code{mutate()} method calls \code{check\_permissions()},
invokes \code{resolve\_mutation()}, and catches
\code{PermissionDeniedError} to return an error payload.

The operation-specific hook is supplied through mixin
inheritance and method resolution order. It is not explicitly
declared as an abstract method in \code{MutationMixin}.

The create and update mixins provide different implementations
of the operation-specific behavior. \code{RelayMutationMixin}
has a similar structure but is not counted as a separate
instance in the inspected inventory.

\subsection{Instance 3: DRF Role Permission Policy}

\textbf{Pattern:} Strategy

\textbf{Category:} Behavioral

\textbf{Origin:} Framework-provided interface with
application-defined policies
\begin{figure}[h]
    \centering
    \includegraphics[width=0.7\textwidth]{uml3.png}
    \caption{Instance 3: DRF Role Permission Policy}
    \label{fig:architecture}
\end{figure}

\textbf{Participants:}

\begin{itemize}[leftmargin=*]
    \item Strategy interface: DRF's
    \code{rest\_framework.permissions.BasePermission}.
    \item Concrete strategies: \code{IsEmployer},
    \code{IsEmployee}, and \code{IsJobCreator}, in
    \code{jobsapp/api/permissions.py:1--26}.
    \item Context and clients: DRF views configuring
    \code{permission\_classes} in the employer and employee
    view modules.
\end{itemize}

\textbf{Intent and evidence:}

REST endpoints require role-based and ownership-based access
rules. The concrete policies implement
\code{has\_permission(request, view)} and are selected by
the relevant views.

The \code{IsJobCreator} policy checks whether the current
user owns the job identified by \code{job\_id}. DRF supplies
the common permission interface and execution mechanism;
the concrete permission rules are application code.

This is a separate instance from the GraphQL permission
hierarchy, even though some policy names overlap.

\subsection{Instance 4: DRF Create Hook}

\textbf{Pattern:} Template Method

\textbf{Category:} Behavioral

\textbf{Origin:} Framework workflow customized by application code

\begin{figure}[h]
    \centering
    \includegraphics[width=0.7\textwidth]{uml4.png}
    \caption{Instance 4: DRF Create Hook}
    \label{fig:architecture}
\end{figure}
\textbf{Participants:}

\begin{itemize}[leftmargin=*]
    \item Framework workflow: DRF's
    \code{CreateModelMixin.create()}.
    \item Hook: \code{perform\_create(serializer)}.
    \item Application overrides:
    \code{JobCreateAPIView.perform\_create()} and
    \code{EmployerCompanyListCreateAPIView.perform\_create()},
    in \code{jobsapp/api/views/employer.py:30--52}.
    \item Additional customization:
    \code{ApplyJobApiView.perform\_create()} and
    \code{ApplyJobApiView.create()}, in
    \code{jobsapp/api/views/employee.py:12--25}.
\end{itemize}

\textbf{Intent and evidence:}

DRF provides a standard resource-creation workflow and
allows application views to customize the save operation
through \code{perform\_create()}.

The employer views attach the current user when saving.
\code{ApplyJobApiView} also overrides \code{create()} and
calls \code{perform\_create()}. It therefore customizes
both the main workflow and the hook.

This is documented as a framework-provided Template Method
workflow with application-specific overrides.

\subsection{Instance 5: Company Serializer Selection}

\textbf{Pattern:} Factory Method candidate

\textbf{Category:} Creational

\textbf{Origin:} Framework selection hook customized by application code
\begin{figure}[h]
    \centering
    \includegraphics[width=0.7\textwidth]{uml5.png}
    \caption{Instance 5: Company Serializer Selection}
    \label{fig:architecture}
\end{figure}

\textbf{Participants:}

\begin{itemize}[leftmargin=*]
    \item Framework creator: DRF's
    \code{GenericAPIView.get\_serializer()}.
    \item Selection hook: \code{get\_serializer\_class()}.
    \item Concrete views: \code{EmployerCompanyListCreateAPIView}
    and \code{EmployerCompanyDetailAPIView}.
    \item Products: \code{CompanySerializer} and
    \code{CompanyWriteSerializer}, both based on
    \code{ModelSerializer}.
\end{itemize}

\textbf{Source locations:}

\code{jobsapp/api/views/employer.py:38--44,55--61} and
\code{jobsapp/api/serializers.py:11--58}.

\textbf{Intent and evidence:}

Company endpoints need different serializers for reading
and writing. The list/create view selects
\code{CompanyWriteSerializer} for POST, while the detail
view selects it for PUT and PATCH. Other methods select
\code{CompanySerializer}.

The selection hook returns a serializer class, which the
framework subsequently instantiates through
\code{get\_serializer()}.

This is retained as a qualified Factory Method candidate
implemented through a framework hook, rather than an
unqualified textbook Factory Method implementation.

\subsection{Rejected and Uncertain Candidates}

The baseline review notes include several candidates that
should not automatically be counted as confirmed patterns:

\begin{itemize}[leftmargin=*]
    \item Applicant status represented by integers and
    conditional branches: conditional logic alone does not
    establish a State pattern.
    \item A status \code{TextChoices} enum by itself:
    an enumeration alone does not establish a behavioral
    State implementation.
    \item \code{RelayMutationMixin}: similar workflow shape,
    but not counted as a separate instance in the inspected
    source scope.
    \item \code{SingleObjectMixin.get\_object}: requires
    further evidence before being classified as a separate
    pattern instance.
\end{itemize}

The inventory above deliberately reports only the five documented
instances. A complete total for rejected or uncertain candidates is not
claimed because the available evidence does not include a reconciled
candidate log.

% ============================================================
\section{Part 2: Evolution and Refactoring}
% ============================================================

Part 2 documents three changes to the Django Job Portal.
Each change is evaluated by its requirement, implementation,
pattern effect, and verification evidence.

\subsection{CH01: Observer Notification}
\begin{figure}
    \centering
    \includegraphics[width=0.7\textwidth]{Observer.png}
    \caption{CH01: Observer Notification}
    \label{fig:architecture}
\end{figure}

\textbf{Requirement:}

Separate job-creation notification behavior from the view
so that notification responsibilities can be extended
without adding each notification action directly to the
job-creation logic.

\textbf{Implementation:}

The recorded change introduces
\code{jobsapp/observers.py}, containing a \code{JobObserver}
protocol/interface, a logging observer, and a subject that
notifies registered observers. The observer behavior is
connected to job creation through the relevant view's
creation hook. Tests are provided in
\code{test\_observers.py}.

\textbf{Pattern effect:}

Observer is recorded as a newly introduced pattern
collaboration, subject to confirmation against the original
baseline source.

\textbf{Design impact:}

The notification responsibility is separated from the
job-creation operation. Additional observers can be
introduced without embedding all notification behavior
in the view.

The trade-offs include registration complexity, additional
indirection, and the need to manage notification failures.
The project notes describe handling observer failures so
that notification failure does not undo a job that has
already been persisted.

\textbf{Verification:}

The supplied CH01 log records the command
\code{python manage.py test tests.jobsapp.test\_observers -v 2}
and reports 6 tests run, all passing. The tests cover observer
subscription, duplicate subscription, notification to multiple observers,
unsubscription, safe unsubscription of an unknown observer, and isolation
of observer failures. The log is retained as
\code{CH01\_Observer\_test\_result.txt}. The before/after commit IDs were
not present in the supplied evidence and are therefore not asserted here.

\subsection{CH02: Recruiter Authorization and GraphQL Query}


\textbf{Requirement:}

Introduce recruiter-specific access control and a GraphQL
query that returns the jobs appropriate to the recruiter.

\textbf{Implementation:}

The recorded change adds an \code{IsRecruiter} permission
candidate and a \code{recruiter\_jobs} query/resolver in
the GraphQL implementation. The intended query restricts
results to published, unfilled jobs.

\begin{figure}[h]
    \centering
    \includegraphics[width=0.7\textwidth]{Strategy.png}
    \caption{CH02: Recruiter Authorization and GraphQL Query}
    \label{fig:architecture}
\end{figure}
\textbf{Pattern effect:}

The permission behavior may be classified as an extension
of the Strategy collaboration
if the source diff confirms
that the policy participates in the same configurable
permission mechanism.


The resolver is recorded as directly instantiating
\code{IsRecruiter}. This is a partial Strategy candidate;
it should not be overstated as a full interchangeable
Strategy family without supporting implementation evidence.


\textbf{Design impact:}

The change makes recruiter authorization explicit and
provides a dedicated query for recruiter jobs. Tests should
cover permitted access, denied access, and the filtering
of returned jobs.


\textbf{Verification:}

The supplied permission test log records
\code{python manage.py test tests.jobsapp.test\_graphql\_permissions -v 2}:
5 tests ran and all passed. They cover anonymous users, authenticated
recruiters, employees, employers, and a missing user. The evidence file is
\code{CH02\_permission\_test\_result.txt}. This log verifies the permission
class, but by itself does not prove end-to-end behavior of the
\code{recruiter\_jobs} GraphQL query or its result filtering. A separate
query integration result should only be reported if its log is attached.

\subsection
{CH03:State for Applicant Status} 
\begin{figure}[h]
    \centering
    \includegraphics[width=0.4\textwidth]{State.png}
    \caption{CH03: State for Applicant Status}
    \label{fig:architecture}
\end{figure}

\textbf{Requirement:}


Make applicant-status transition rules explicit and
validate changes to application status.

\textbf{Implementation:}

The recorded change introduces
\code{jobsapp/application\_states.py}, including Pending,
Accepted, and Rejected state classes and an
\code{ApplicationStateFactory}. The related endpoint is
\code{UpdateApplicantStatusAPIView}.


\textbf{Pattern effect:}

State is recorded as a newly introduced pattern candidate.
The classification depends on state objects actually
encapsulating state-dependent behavior and being used by
the application, rather than merely representing status
labels or returning state objects from a factory.

\textbf{Design impact:}

Explicit state objects can make permitted transitions
easier to understand and test. The trade-off is additional
classes and indirection, which may be unnecessary for a
very small lifecycle.

Tests should cover valid transitions, invalid transitions,
terminal-state behavior, and unknown status codes.



\textbf{Verification:}

The supplied CH03 log reports 13 tests run, all passing, covering the
application-state tests and related employer API-view tests. Coverage
includes Pending-to-Accepted and Pending-to-Rejected transitions,
terminal behavior for Rejected, Accepted-to-Rejected behavior, status
codes and labels, unknown status codes, and API authentication/not-found
handling. The evidence file is \code{CH03\_State\_test\_result.txt}.
These tests support the transition behavior implemented for this change;
the State classification is retained because the change introduces
Pending, Accepted, and Rejected state classes and an
\code{ApplicationStateFactory}.


\subsection{Pattern-Count Evolution}

The project pattern-count CSV records the following
evolution across the original baseline and three changes.

\begin{center}
\begin{tabularx}{\linewidth}{@{}lrrrrrrr@{}}
\toprule
Snapshot & Factory & Obs. & State & Strat. & T.M. & Types & Inst.\\
\midrule
Original & 1 & 0 & 0 & 2 & 2 & 3 & 5\\
After CH01 & 1 & 1 & 0 & 2 & 2 & 4 & 6\\
After CH02 & 1 & 1 & 0 & 3 & 2 & 4 & 7\\
After CH03 & 1 & 1 & 1 & 3 & 2 & 5 & 8\\
\bottomrule
\end{tabularx}
\end{center}

\textbf{Abbreviations:} Obs. = Observer; Strat. = Strategy;
T.M. = Template Method; Inst. = instances.

These are the counts recorded in the project CSV. The
Factory Method classification remains qualified, and the
CH02 Strategy entry should retain its partial-candidate
qualification unless the source review establishes the
stronger classification.

The instance count should represent independent
collaborations, not simply the number of classes added.
For each snapshot, the final inventory should preserve the
name of a continuing collaboration and explain whether
it was preserved, modified, extended, replaced, removed,
or newly introduced.

\subsection{Before-and-After Metrics}

Test results below come from the supplied logs or previously recorded
project notes. Source LOC values were not included in the available
artifacts, so they are marked as not measured rather than estimated.

\begin{longtable}{@{}p{0.27\linewidth}p{0.22\linewidth}p{0.22\linewidth}p{0.21\linewidth}@{}}
\toprule
Snapshot & Source LOC & Test result & Evidence\\
\midrule
Original baseline &
Not measured &
31 passed, 5 failed, 5 errors among 41 executed; 18 not executed &
Baseline test log\\

After CH01 &
Not measured &
6 passed, 0 failed &
\code{CH01\_Observer\_test\_result.txt}\\

After CH02 &
Not measured &
5 permission tests passed, 0 failed; query integration not confirmed by this log &
\code{CH02\_permission\_test\_result.txt}\\

After CH03 &
Not measured &
13 passed, 0 failed in the supplied state/API test run &
\code{CH03\_State\_test\_result.txt}\\
\bottomrule
\end{longtable}

The baseline run also reported 18 tests not executed because
four test classes failed during \code{setUpClass}. The
separate Part 1 pattern-test suite reported 17 passing tests.
These results refer to different test runs and must not be
combined into a single test total.

The source-LOC comparison remains unreported because matching source
snapshots and measured counts were not supplied. These values should be
calculated from the baseline and each post-change revision using the same
file scope before making quantitative claims about code growth.

\subsection{Overall Interpretation}

The three changes aim to improve separation of concerns,
authorization clarity, and explicit lifecycle behavior.

Observer separates notification handling from job creation.
The recruiter query change makes a specific access rule and
query behavior explicit. State models applicant-status
transitions through dedicated state objects.

These structures introduce additional classes and indirection.
Whether that cost is justified depends on the complexity of
the requirements and the evidence from tests and maintenance.
An increase or decrease in source LOC alone does not prove
that a design is better.


% ============================================================

% ============================================================
\section{Part 3: Cross-Language Implementations}
% ============================================================

\subsection{Objective and scope}

Part 3 implements three selected patterns independently in Python, Java,
JavaScript, and C++ so that their structures and measured characteristics
can be compared across languages. The selected patterns are Observer
(based on the Part 2 job-creation notification), Strategy (based on the
DRF role-permission policy), and Template Method (based on the GraphQL
mutation pipeline). The Part 3 prompt explicitly excludes State from this
cross-language implementation scope.

These are small, standalone reproductions intended to demonstrate pattern
structure and behavior. They are not complete integrations with Django,
Django REST Framework, or Graphene. The Strategy implementation is inspired
by DRF permission behavior; the Observer implementation models job-created
notifications. The original GraphQL mutation pipeline checks permissions,
calls \code{resolve\_mutation()} when allowed, and converts
\code{PermissionDeniedError} into an error payload.

\subsection{Implementation overview}

\begin{longtable}{p{2.2cm}p{3.0cm}p{8.0cm}}
\toprule
\textbf{Pattern} & \textbf{Reference behavior} & \textbf{Cross-language structure} \\
\midrule
\endfirsthead
\toprule
\textbf{Pattern} & \textbf{Reference behavior} & \textbf{Cross-language structure} \\
\midrule
\endhead
Observer & Job creation triggers notification & A job service maintains registered observers and notifies them when a job is created. Email-style and analytics observers illustrate separate reactions to the same event. \\
\addlinespace
Strategy & DRF role and ownership permissions & A common permission-strategy interface is implemented by employer, employee, and job-creator checks. The caller delegates the authorization decision to the selected strategy. \\
\addlinespace
Template Method & GraphQL mutation pipeline & A base pipeline fixes the sequence of processing while subclasses customize the mutation hook. The Java and JavaScript versions describe permission-check, hook execution, and permission-denial payload behavior; the current Python and C++ examples are generic pipelines and are not yet faithful equivalents of the original flow. \\
\bottomrule
\end{longtable}

\subsection{Source and test inventory}

The submitted archive contains implementation files, test runners, and a README in each language folder for the three patterns. The metrics below are transcribed from the supplied \code{cross\_language.csv} and \code{llm.csv} records and language READMEs. LOC means the recorded source lines for the implementation file; the supplied records do not document a consistent blank-line/comment-line counting rule, so LOC comparisons should be treated as provisional.

\begin{longtable}{p{2.0cm}p{1.6cm}p{1.6cm}p{1.8cm}p{2.5cm}p{3.6cm}}
\toprule
\textbf{Pattern} & \textbf{Language} & \textbf{Source file} & \textbf{LOC} & \textbf{Test result recorded} & \textbf{Important qualification} \\
\midrule
\endfirsthead
\toprule
\textbf{Pattern} & \textbf{Language} & \textbf{Source file} & \textbf{LOC} & \textbf{Test result recorded} & \textbf{Important qualification} \\
\midrule
\endhead
Observer & Python & \code{main.py} & 27 & 3/3 passed (recorded) & Standalone reproduction. \\
Observer & Java & \code{Main.java} & 40 & 3/3 passed (recorded) & JVM startup included in timing; compilation excluded. \\
Observer & JavaScript & \code{main.js} & 34 & 3/3 passed (recorded) & Node.js startup included in timing. \\
Observer & C++ & \code{main.cpp} & 47 & 3/3 passed (recorded) & Maximum function CCN recorded as 2 by Lizard 1.24.3. \\
\addlinespace
Strategy & Python & \code{main.py} & 27 & 3/3 passed (recorded) & Standalone; not a complete DRF integration. \\
Strategy & Java & \code{Main.java} & 46 & 3/3 passed (recorded) & Standalone; not a complete DRF integration. \\
Strategy & JavaScript & \code{main.js} & 31 & 3/3 passed (recorded) & Standalone; not a complete DRF integration. \\
Strategy & C++ & \code{main.cpp} & 61 & \textbf3/3 passed (recorded) &  Standalone; not a complete DRF integration. \\
\addlinespace
Template Method & Python & \code{main.py} & 33 & 3/3 passed (recorded) & Generic validation pipeline; not faithful to original GraphQL flow yet. \\
Template Method & Java & \code{Main.java} & 40 & 3/3 passed (recorded) & Standalone reproduction of permission-check, mutation hook, and denial payload. \\
Template Method & JavaScript & \code{main.js} & 38 & 3/3 passed (recorded) & Standalone reproduction of permission-check, mutation hook, and denial payload. \\
Template Method & C++ & \code{main.cpp} & 76 & 3/3 passed (recorded) & Generic pipeline; not faithful to original GraphQL flow yet; maximum function CCN recorded as 3. \\
\bottomrule
\end{longtable}

\subsection{Recorded runtime and toolchain information}

The supplied CSV records the following runtime/compiler versions and median
process times. Times are in milliseconds and are the recorded medians, not
independent performance guarantees. Python and Node.js timings include
process startup. Java timings include JVM startup and exclude compilation.
C++ timings include process startup and execution, with compilation excluded.
The archive does not include the raw five-run timing series for every row,
so those individual observations cannot be independently checked here.

\begin{longtable}{p{2.2cm}p{2.6cm}p{3.6cm}p{2.4cm}p{3.2cm}}
\toprule
\textbf{Pattern} & \textbf{Language} & \textbf{Recorded version} & \textbf{Median (ms)} & \textbf{Complexity record} \\
\midrule
\endfirsthead
\toprule
\textbf{Pattern} & \textbf{Language} & \textbf{Recorded version} & \textbf{Median (ms)} & \textbf{Complexity record} \\
\midrule
\endhead
Observer & Python & Python 3.13.15 & 64.955 &  Max CCN 2 (Lizard 1.24.3)\\
Observer & Java & OpenJDK 25.0.5.0.2 LTS & 98.713 &  Max CCN 2 (Lizard 1.24.3)\\
Observer & JavaScript & Node.js v24.22.1 & 82.355 &  Max CCN 2 (Lizard 1.24.3)\\
Observer & C++ & GCC 16.2.0 (MSYS2) & 102.5902 & Max CCN 2 (Lizard 1.24.3) \\
\addlinespace
Strategy & Python & Python 3.13.15 & 66.214 &Max CCN 3 (Lizard 1.24.3) \\
Strategy & Java & OpenJDK 25.0.5.0.2 LTS & 83.283 & Max CCN 3 (Lizard 1.24.3) \\
Strategy & JavaScript & Node.js v24.22.1 & 108.211 & Max CCN 6 (Lizard 1.24.3) \\
Strategy & C++ & GCC 16.2.0 (MSYS2) & 104.2196 & Max CCN 3 (Lizard 1.24.3) \\
\addlinespace
Template Method & Python & Python 3.13.15 & 88.082 & Max CCN 3 (Lizard 1.24.3) \\
Template Method & Java & OpenJDK 25.0.5.0.2 LTS & 73.920 & Max CCN 2 (Lizard 1.24.3) \\
Template Method & JavaScript & Node.js v24.22.1 & 89.399 & Max CCN 3 (Lizard 1.24.3) \\
Template Method & C++ & GCC 16.2.0 (MSYS2) & 114.1219 & Max CCN 3 (Lizard 1.24.3) \\
\bottomrule
\end{longtable}

\subsection{Cross-language observations}

\begin{itemize}[leftmargin=*]
    \item \textbf{Python:} Concise class and method definitions make the pattern roles easy to inspect. The recorded implementation LOC is low in the Observer and Strategy examples.
    \item \textbf{Java:} Explicit interfaces and class declarations make the pattern contracts visible. The recorded runtime measurements include JVM startup, which affects very small programs.
    \item \textbf{JavaScript:} The examples use JavaScript classes and Node.js execution. The recorded timings include Node.js process startup.
    \item \textbf{C++:} Abstract base classes and virtual methods demonstrate runtime polymorphism. Lizard complexity values are present for the three C++ implementations, but comparable complexity measurements are not recorded for all other languages.
    \item \textbf{Timing caution:} The programs are small, and the reported medians include runtime startup. These measurements should not be used to conclude that one language or pattern is intrinsically faster without matching workloads, repeated raw measurements, and controlled environments.
\end{itemize}

\subsection{Part 3 evidence gaps and follow-up corrections}

The Part 3 artifacts identify several items that should remain explicit in the report rather than being presented as complete:

\begin{enumerate}[leftmargin=*]
    \item The Strategy C++ row in \code{cross\_language.csv} says ``VERIFY TEST RUNNER''. Run its documented test command and update the CSV and README with the actual result before submission.
    \item The Template Method Python implementation uses a generic validation pipeline, and the C++ implementation is also labelled generic. To match the reference pattern, both should implement the permission-check $\rightarrow$ \code{resolve\_mutation()} hook $\rightarrow$ permission-denial error-payload sequence.
    \item The supplied \code{cross\_language.csv} has malformed rows where commas in unquoted notes merge adjacent records. The metrics table above uses clearly recoverable values and the accompanying LLM log/README records; repair the CSV quoting before using it as a machine-readable final dataset.
    \item \code{pattern\_crosslanguage.csv} contains only its header in the supplied archive. Populate it from verified measurements or clearly designate \code{cross\_language.csv} as the authoritative comparison file.
    \item The prompt requests five local timing runs, raw measurements, consistent LOC rules, and actual version output. The archive records medians and version strings but does not preserve all raw timing observations or a clear LOC counting rule. These should be added if required for reproducibility.
    \item The LLM provenance log explicitly says that exact historical prompts were not captured and that task summaries were reconstructed. This limitation is preserved; historical prompt text has not been invented.
\end{enumerate}

\subsection{Part 3 conclusion}

Part 3 provides cross-language examples for Observer, Strategy, and Template Method in Python, Java, JavaScript, and C++. The examples illustrate common pattern roles while also showing language-specific expression through interfaces, classes, virtual dispatch, and runtime conventions. The current evidence supports a descriptive comparison, but not a definitive performance ranking: one Strategy C++ test result remains unverified, and the Python and C++ Template Method examples still require alignment with the original GraphQL mutation pipeline. After those corrections and consistent metric collection, the comparison can serve as a reproducible extension of the Part 1 identification and Part 2 pattern evolution.


\appendix
\section{LLM Records and Human Verification}

Free and publicly accessible tools were used to assist with
design pattern analysis, implementation, and testing throughout
the case study. LLM-generated suggestions were treated as
proposals and were reviewed against the source code and test
results before being accepted.

\subsubsection{Model and Execution Records}

The case study was divided into three parts. Part 1 examined
design patterns in the original Django Job Portal repository.
Part 2 introduced and evaluated changes involving the Observer
pattern for job-creation notifications, recruiter authorization
in GraphQL, and the State pattern for application status
transitions. Part 3 implemented standalone examples of selected
design patterns in Python, Java, JavaScript, and C++.

For reproducibility, each major LLM-assisted analysis should
be associated with its displayed model, access tier, execution
date, exact prompt, supplied input, source file paths, and
resulting output. The original repository commit used for the
baseline was
\texttt{f6e5efb33e9aac7334a364ee46d6e2ad1321db30}.
Where historical prompts or raw outputs were not preserved,
they are recorded as unavailable rather than reconstructed
as original experimental evidence.

\subsubsection{Human Verification of Design Pattern Claims}

The identified patterns and proposed implementations were
evaluated by examining their structure, responsibilities,
collaborations, and behavior in the code.

For Part 1, pattern claims were compared with the original
repository implementation. For Part 2, the Observer
implementation was examined in
\texttt{jobsapp/observers.py}, the authorization-related
GraphQL changes in
\texttt{jobsapp/graphql/permissions.py} and
\texttt{jobsapp/graphql/queries.py}, and the State
implementation in \texttt{jobsapp/application\_states.py}.
The corresponding focused tests were used as behavioral
evidence. The GraphQL authorization changes were evaluated
as permission enforcement and query integration; they were
not automatically treated as proof of a complete Strategy
pattern implementation.

Claims were assessed using four categories: correct,
partially correct, incorrect, and cannot confirm. A claim
was considered correct when supported by the implementation
and available verification evidence. Partially correct
claims required qualification or correction, while claims
without sufficient evidence were classified as cannot
confirm. Human review was used to identify overgeneralized
pattern classifications and distinguish structural evidence
from successful test execution.

\subsubsection{Generated Code and Test Evidence}

LLM-assisted code was reviewed and tested using the project's
available test suite. The recorded focused test results were
as follows:

\begin{itemize}
    \item \textbf{Observer:} All 6 tests passed, covering
    subscription, notification, unsubscription, duplicate
    subscriptions, and observer failure handling.

    \item \textbf{GraphQL recruiter query:} All 4 tests passed,
    covering anonymous access, employer access, the existing
    public jobs query, and recruiter access to published,
    unfilled jobs.

    \item \textbf{GraphQL permissions:} All 5 tests passed,
    covering anonymous users, authenticated recruiters,
    employees, employers, and missing users.

    \item \textbf{State and API tests:} The recorded focused
    run passed all 13 tests, including the five State-specific
    tests for valid transitions, terminal rejection behavior,
    status labels, and unknown status codes.
\end{itemize}

A separate combined State and API test run executed 20 tests.
Nineteen passed and one failed:
\texttt{test\_jobs\_list\_api\_view}. The failure occurred
because the endpoint returned four jobs while the test
expected five. This discrepancy was retained as an unresolved
test issue rather than reported as a successful result.
The passing pattern-specific tests therefore do not establish
that the entire project's test suite passes.

\subsubsection{Reproducibility and Limitations}

The supporting records include focused test outputs and
documentation associated with the three Part 2 changes.
For complete reproducibility, each major analysis should
also preserve its exact prompt, supplied source context,
unaltered LLM output, model and access-tier information,
execution date, reviewer assessment, and any correction made.

The case study distinguishes recorded evidence from
retrospective explanations. Missing experimental details
are not presented as though they had been captured at the
time of execution. Test outcomes are reported according to
the actual recorded results, and unresolved failures are
retained as limitations for further investigation.

% ============================================================
\section{Conclusion}
% ============================================================

The documented baseline inventory contains five instances
across three pattern types: Strategy, Template Method, and
a qualified Factory Method candidate.

The two baseline Strategy instances are the GraphQL Permission Policy
and the DRF Role Permission Policy. CH02 adds a recruiter permission
that is treated as a partial Strategy candidate rather than a confirmed
third full Strategy collaboration. The two Template Method
instances are the GraphQL Mutation Pipeline and the DRF
Create Hook. Company Serializer Selection is retained as
a qualified Factory Method candidate because it relies on
a DRF selection hook.

Part 2 records three changes: Observer-based job-creation
notifications, recruiter-specific authorization/query
behavior, and state-based applicant-status handling. The
recorded pattern-count CSV shows an evolution from three
types and five instances in the baseline to five types and
eight instances after CH03. The qualifications noted in
the pattern analysis must be preserved when interpreting
these counts.

The strength of the case study rests on the traceability
of each pattern to source code, the justification of its
classification, and the verification of behavior through
tests and change evidence.

% ============================================================
\section{Contributions and References}
% ============================================================

\subsection{Team Contributions}

\begin{tabularx}{\linewidth}{@{}lXl@{}}
\toprule
Roll Number & Student Name & Contribution\\
\midrule
AM.SC.U4CSE23233 & M Poojith & Report\\
AM.SC.U4CSE23234 & M Surya Prakash & Part1\\
AM.SC.U4CSE23241 & P Rushikesh & Part2\\
AM.SC.U4CSE23247 & S Krishna Vamsi & Part3\\
\bottomrule
\end{tabularx}

\subsection{References}

\begin{enumerate}[leftmargin=*]
    \item Django REST Framework, Permissions:
    \url{https://www.django-rest-framework.org/api-guide/permissions/}

    \item Django REST Framework, Generic Views:
    \url{https://www.django-rest-framework.org/api-guide/generic-views/}

    \item Graphene-Python documentation:
    \url{https://docs.graphene-python.org/}

    \item Case study repository:
    \url{https://github.com/aswathymohan-amrita/23CSE455_DP_CaseStudy}

    \item Gamma, E., Helm, R., Johnson, R., and Vlissides, J.,
    \textit{Design Patterns: Elements of Reusable Object-Oriented Software},
    Addison-Wesley, 1994. Confirm the edition required by the course.
\end{enumerate}


\end{document}